\chapter{Due paesaggi, un contrasto}
\label{cap:due_paesaggi}

\section{Una regione, non un singolo sito}

Chi ha viaggiato nel deserto di Atacama e nel Cile settentrionale porta con sé un'impressione difficile da restituire a parole: l'arte rupestre, in quel territorio, non è un monumento isolato da visitare come si visita una cattedrale. È un fenomeno diffuso, presente in almeno sei contesti ambientali molto diversi tra loro, ciascuno con la propria storia, la propria tecnica, i propri motivi ricorrenti. Vale la pena percorrerli uno per uno, perché è proprio questa varietà di contesti, più ancora della quantità assoluta di segni, a rendere il confronto con le Alpi venete istruttivo.

\subsection{Lungo le piste di quota: il Qhapaq Ñan}

Il Camino del Inca, o Qhapaq Ñan, attraversava longitudinalmente anche il nord del Cile, con un ramo che dalla costa saliva verso le terre alte dell'Atacama e un ramo più interno, più lungo e più frequentato, che passava per l'attuale Bolivia e il nord-ovest argentino. Lungo questo tracciato e i suoi rami secondari sorgevano i tambos, punti di rifornimento e sosta dove i viaggiatori potevano riposare e dove le autorità imperiali amministravano i territori vicini.

Solo nel tratto cileno della rete, i cui resti sono oggi riconosciuti patrimonio dell'umanità dall'UNESCO, sono stati censiti 112 chilometri di percorso e 138 siti archeologici associati \citep{mnhn_qhapaq_atacama}. Lungo un tratto che risale da Lasana verso nord, appena a settentrione della Región de Atacama, si incontrano più siti di arte rupestre in successione, tra cui quelli noti come Santa Bárbara, Incahuasi e Cerro Colorado: segno che i tracciati stessi del cammino imperiale, e i tratti di pista preesistenti che l'impero inca aveva incorporato, erano costellati di segni incisi o dipinti, non solo nei punti di sosta principali.

\subsection{Lungo la costa}

Un secondo contesto è quello costiero. Presso Caldera, vicino alla caletta di Obispo, si trova il sito di Las Lizas, studiato già negli anni ottanta da Hans Niemeyer \citep{niemeyer1985}: un insieme di pannelli incisi con raffigurazioni di specie marine, associati a resti di conchiglie e a materiale in pietra e ceramica. È un caso interessante perché mostra come l'arte rupestre della regione non fosse legata solo alla pastorizia di quota o al commercio carovaniero, ma anche a comunità che vivevano di risorse marine, con un repertorio di soggetti (pesci, probabilmente cetacei o altre specie marine) del tutto diverso da quello dei camelidi e delle figure umane più tipico delle zone interne.

Più a nord, nella Cordillera di Taltal (Región de Antofagasta), il sito di El Médano presenta un caso ancora più straordinario per estensione: pitture rupestri che si susseguono senza soluzione di continuità per circa cinque chilometri lungo il letto di una quebrada, per un dislivello di circa quattrocento metri, sempre al di sopra del limite della nebbia e della vegetazione costiera \citep{museo_precolombino_medano}. Non un singolo pannello, ma una sequenza quasi ininterrotta di immagini lungo l'intero percorso di un corso d'acqua stagionale.

\subsection{Nelle antiche miniere}

Un terzo contesto, forse il meno intuitivo per chi pensa all'arte rupestre solo in termini di percorsi o di insediamenti, è quello minerario. Nel deserto di Atacama, comunità di minatori di piccola scala, attive già prima dell'arrivo degli Inca, estraevano materiali come la malachite e la turchesa per la produzione di collane, e pigmenti rossi per uso rituale o decorativo. Quando l'impero inca costruì la propria rete stradale attraverso questi stessi territori, queste comunità minerarie non furono semplicemente assorbite: riorganizzarono i propri insediamenti produttivi proprio in funzione della vicinanza al nuovo tracciato, per facilitare lo scambio dei propri manufatti, installando campi con architettura locale direttamente lungo il Qhapaq Ñan \citep{qhapaq_atacama_minas}. È un caso che mostra come i luoghi di lavoro, non solo i luoghi di passaggio o di sosta, fossero segnati dalla stessa pratica di incidere e dipingere la roccia.

Va notato che l'area di Taltal, già menzionata per il sito di El Médano, ha restituito anche alcune delle prove più antiche di attività mineraria umana finora conosciute al mondo. Tra le colline che scendono verso la costa, in un raggio di pochi chilometri da Taltal, sono state identificate cavità di estrazione di ocra rossa, uno dei pigmenti usati nella pittura rupestre e in contesti rituali, che alcune datazioni radiocarboniche collocano tra 10.000 e 12.000 anni prima dell'anno zero, con possibili tracce di attività ancora più antiche \citep{taltal_minas_antiguas}. Si tratta di un dato che rovescia una prospettiva intuitiva: non solo l'arte rupestre di questa costa è straordinariamente diffusa, ma il territorio circostante ha prodotto anche le miniere più antiche oggi note, e lo ha fatto in un periodo cronologicamente paragonabile, almeno nella sua fase finale, all'inizio della frequentazione mesolitica delle Alpi venete. I due fenomeni non sono direttamente collegati, e nessuno sostiene che vi sia una relazione causale; ma la coincidenza temporale è utile per ricordare che il confronto tra Atacama e Alpi non è un confronto tra un'area ``ricca di storia'' e una ``priva di storia'', bensì tra due storie ugualmente profonde, di cui una è stata molto più intensamente cercata e documentata dell'altra.

\subsection{Nel deserto interno}

Un quarto contesto, infine, è quello più propriamente definibile come deserto interno, lontano sia dalla costa sia dai grandi tracciati principali: pitture e petroglifi sparsi lungo piste minori, oggi meno note e meno frequentate di un tempo, che collegavano oasi, sorgenti e punti di passaggio secondari. Nella parte settentrionale della Región de Atacama, tra il margine meridionale del cosiddetto Despoblado e il fiume Copiapó, le pitture rupestri sono particolarmente frequenti e vengono collegate al complesso culturale Las Ánimas, attivo tra il settecento e il milleduecento dopo l'anno zero \citep{memoriachilena_rupestre}. Più a monte, nell'alta valle del fiume Huasco, i petroglifi si concentrano con motivi animali, umani e soprattutto geometrici, in parte riferibili allo stile denominato La Silla; frequenti, in questa zona, sono anche piccole coppelle circolari, chiamate localmente ``copitas'', spesso raggruppate sulla parete verticale della roccia.

\subsection{Nelle valli interne: tre casi meno noti}

Il quinto e il sesto contesto appartengono ugualmente al deserto interno, ma meritano una menzione separata perché mostrano come l'arte rupestre di quest'area compaia anche in luoghi non riconducibili a nessuno dei contesti già descritti.

Nella Quebrada de las Pinturas, lungo un affluente minore del sistema del Copiapó, le pitture rupestri comprendono una concentrazione insolita di figure note come ``sacrificadores'': rappresentazioni umane in postura rituale, probabilmente legate a cerimonie il cui significato è ancora discusso. È un repertorio iconografico distinto da quello dei camelidi e delle scene di caccia dominanti in altre zone, che mostra come l'arte rupestre di quest'area non risponda a un unico codice tematico.

Nei pressi del Llano de Varas, il sito di Cachiyuyo de Oro ha restituito pitture su pareti rocciose riparate, in un ambiente di pianura e bassa collina che differisce nettamente dai siti di alta quota o costieri già descritti. Qualche distanza più a nord, nell'area di Chañaral Alto, la proprietà agricola nota come Finca de Chañaral Alto conserva pitture policromate su rocce affioranti in un ambiente di fondo valle, in un contesto dove l'arte rupestre è intrecciata con i segni delle comunità di agricoltori e allevatori che hanno occupato queste valli per millenni. Insieme, questi tre siti mostrano che il fenomeno dell'arte rupestre in quest'area non conosce un unico tipo di paesaggio: si manifesta lungo i grandi cammini di quota, sulle coste, nelle miniere, nel deserto interno, nelle valli e persino nelle proprietà agrarie che ancora oggi occupano i fondivalle.

\section{Lo stesso tipo di ambiente, un risultato diverso}

Chi si sposta poi sulle Alpi venete, in alcuni contesti montani comparabili per funzione, come valichi, ripari e affioramenti rocciosi, fa fatica a trovare qualcosa di simile a uno solo di questi contesti, figuriamoci a tutti e quattro insieme. Non che manchi del tutto, come mostrerà il terzo capitolo di questo libro: esistono alcuni luoghi noti, alcuni nomi che ricorrono nella letteratura locale. Ma il confronto è impietoso non solo per la quantità di segni, ma per la varietà di contesti in cui quei segni compaiono. Nel deserto di Atacama e nel Cile settentrionale l'arte rupestre accompagna le piste di quota, la costa, le miniere, il deserto interno e le valli agricole, ambienti radicalmente diversi tra loro. Nelle Alpi venete, come si vedrà, i pochi nuclei noti si concentrano quasi tutti in un unico tipo di contesto, quello collinare-prealpino intorno al Garda.

Questo libro nasce da quel doppio confronto, di quantità e di varietà, e dalla domanda che inevitabilmente suscita: perché? È possibile che due territori entrambi montani, entrambi frequentati per millenni da cacciatori, pastori e mercanti, abbiano prodotto un'eredità così diversa? O è più probabile che la differenza non sia nel passato, ma nel presente: nella quantità e nella qualità della ricerca che è stata dedicata a cercarla?

\section{Tre variabili, non una sola}

Di fronte a un contrasto così netto, la tentazione è saltare subito a una conclusione unica. È una tentazione da maneggiare con attenzione, perché il contrasto tra la Región de Atacama e le Alpi venete nasconde almeno tre variabili distinte, che vale la pena separare ed esaminare una per una, con le evidenze che le sostengono, invece di lasciarle come ipotesi generiche.

\subsection{La roccia}

La prima variabile è la litologia, cioè il tipo di roccia disponibile. Le superfici usate nei quattro contesti appena descritti sono in gran parte rocce vulcaniche scure, spesso ricoperte da una patina scura di ossidi (la cosiddetta patina del deserto), che un colpo deciso con uno strumento di pietra più dura asporta, esponendo lo strato più chiaro sottostante e creando un contrasto visivo netto e duraturo. È una tecnica che funziona particolarmente bene proprio su questo tipo di superficie, ed è coerente con la tecnica osservata anche nei siti costieri e minerari, non solo in quelli di quota.

Le rocce delle Alpi venete sono in prevalenza calcaree, con alcune aree vulcaniche (i Colli Euganei, discussi più avanti in questo libro) dominate da trachiti e altre vulcaniti compatte, meno inclini a produrre lo stesso contrasto visivo netto quando incise. Non è un ostacolo insormontabile, come dimostrano le incisioni del Garda su roccia calcarea, ma è una differenza reale, che rende la tecnica meno immediatamente efficace rispetto al caso atacameno.

\subsection{La conservazione}

La seconda variabile è la conservazione nel tempo. La Región de Atacama comprende, per ampi tratti, uno degli ambienti più aridi del pianeta: la quasi totale assenza di pioggia rallenta enormemente i processi che cancellano un'incisione, come l'azione del gelo e del disgelo, la crescita di licheni, l'espansione delle radici delle piante. Anche nei siti costieri, dove l'aridità è mitigata dall'umidità marina, il clima resta comunque molto più stabile e meno aggressivo di quello alpino.

Il clima delle Alpi venete lavora nella direzione opposta: precipitazioni abbondanti, cicli di gelo e disgelo frequenti in quota, vegetazione che ricolonizza rapidamente ogni superficie non periodicamente ripulita. Un'incisione poco profonda, come si vedrà a proposito della Val d'Assa nel terzo capitolo, può diventare illeggibile in pochi secoli, anche quando la sua esistenza è accertata.

\subsection{La ricerca}

La terza variabile, quella che interessa di più in questo libro, è la storia della ricerca. Un'incisione rupestre non si offre da sola all'osservatore: va cercata, spesso in luoghi scomodi da raggiungere, con un occhio allenato a riconoscere un segno intenzionale da un graffio naturale della roccia. Se in una regione decine di studiosi hanno passato una vita a percorrere sistematicamente ogni tipo di ambiente, dalle piste di quota alla costa alle aree minerarie, e in un'altra regione lo stesso lavoro non è mai stato fatto con la stessa intensità, il numero di siti noti rifletterà tanto la storia della ricerca quanto la storia della frequentazione umana antica. È il tema a cui è dedicato per intero il prossimo capitolo.

\section{Un indizio concreto, prima di proseguire}

Queste tre variabili non si escludono a vicenda. È perfettamente possibile che tutte e tre contribuiscano al contrasto osservato tra la Región de Atacama e le Alpi venete, in proporzioni che nessuno ha ancora misurato con precisione per questo specifico confronto. Il compito di questo libro non è dimostrare che una sola di esse spiega tutto, ma capire quanto la terza, la più controllabile delle tre, possa aver pesato più di quanto si creda, e cosa si potrebbe fare per verificarlo.

Un indizio a favore di questa ipotesi, va detto fin da subito, non è solo un ragionamento astratto. Come mostrerà con dati concreti il prossimo capitolo, esistono casi documentati in cui il sospetto di un research bias è stato messo alla prova sul campo, con esito netto: aree considerate fino a pochi anni fa marginali e prive di frequentazione preistorica si sono rivelate, dopo una ricerca sistematica dedicata, dense di testimonianze mai prima cercate con la stessa cura. Questi casi non dimostrano che lo stesso accadrà in Veneto. Dimostrano però che la scarsità di siti noti, da sola, non è mai una prova di scarsità reale: è, nella migliore delle ipotesi, un dato ambiguo, che può nascondere tanto un'assenza autentica quanto una ricerca mai svolta con la necessaria intensità.
